% Options for packages loaded elsewhere
\PassOptionsToPackage{unicode}{hyperref}
\PassOptionsToPackage{hyphens}{url}
\PassOptionsToPackage{dvipsnames,svgnames,x11names}{xcolor}
%
\documentclass[
  11pt,
  a4paper,
]{report}

\usepackage{amsmath,amssymb}
\usepackage{setspace}
\usepackage{iftex}
\ifPDFTeX
  \usepackage[T1]{fontenc}
  \usepackage[utf8]{inputenc}
  \usepackage{textcomp} % provide euro and other symbols
\else % if luatex or xetex
  \usepackage{unicode-math}
  \defaultfontfeatures{Scale=MatchLowercase}
  \defaultfontfeatures[\rmfamily]{Ligatures=TeX,Scale=1}
\fi
\usepackage{lmodern}
\ifPDFTeX\else  
    % xetex/luatex font selection
    \setmainfont[]{TeX Gyre Pagella}
    \setsansfont[]{TeX Gyre Heros}
\fi
% Use upquote if available, for straight quotes in verbatim environments
\IfFileExists{upquote.sty}{\usepackage{upquote}}{}
\IfFileExists{microtype.sty}{% use microtype if available
  \usepackage[]{microtype}
  \UseMicrotypeSet[protrusion]{basicmath} % disable protrusion for tt fonts
}{}
\makeatletter
\@ifundefined{KOMAClassName}{% if non-KOMA class
  \IfFileExists{parskip.sty}{%
    \usepackage{parskip}
  }{% else
    \setlength{\parindent}{0pt}
    \setlength{\parskip}{6pt plus 2pt minus 1pt}}
}{% if KOMA class
  \KOMAoptions{parskip=half}}
\makeatother
\usepackage{xcolor}
\usepackage[top=1in,bottom=1in,left=1.25in,right=1.25in]{geometry}
\setlength{\emergencystretch}{3em} % prevent overfull lines
\setcounter{secnumdepth}{5}
% Make \paragraph and \subparagraph free-standing
\makeatletter
\ifx\paragraph\undefined\else
  \let\oldparagraph\paragraph
  \renewcommand{\paragraph}{
    \@ifstar
      \xxxParagraphStar
      \xxxParagraphNoStar
  }
  \newcommand{\xxxParagraphStar}[1]{\oldparagraph*{#1}\mbox{}}
  \newcommand{\xxxParagraphNoStar}[1]{\oldparagraph{#1}\mbox{}}
\fi
\ifx\subparagraph\undefined\else
  \let\oldsubparagraph\subparagraph
  \renewcommand{\subparagraph}{
    \@ifstar
      \xxxSubParagraphStar
      \xxxSubParagraphNoStar
  }
  \newcommand{\xxxSubParagraphStar}[1]{\oldsubparagraph*{#1}\mbox{}}
  \newcommand{\xxxSubParagraphNoStar}[1]{\oldsubparagraph{#1}\mbox{}}
\fi
\makeatother


\providecommand{\tightlist}{%
  \setlength{\itemsep}{0pt}\setlength{\parskip}{0pt}}\usepackage{longtable,booktabs,array}
\usepackage{calc} % for calculating minipage widths
% Correct order of tables after \paragraph or \subparagraph
\usepackage{etoolbox}
\makeatletter
\patchcmd\longtable{\par}{\if@noskipsec\mbox{}\fi\par}{}{}
\makeatother
% Allow footnotes in longtable head/foot
\IfFileExists{footnotehyper.sty}{\usepackage{footnotehyper}}{\usepackage{footnote}}
\makesavenoteenv{longtable}
\usepackage{graphicx}
\makeatletter
\def\maxwidth{\ifdim\Gin@nat@width>\linewidth\linewidth\else\Gin@nat@width\fi}
\def\maxheight{\ifdim\Gin@nat@height>\textheight\textheight\else\Gin@nat@height\fi}
\makeatother
% Scale images if necessary, so that they will not overflow the page
% margins by default, and it is still possible to overwrite the defaults
% using explicit options in \includegraphics[width, height, ...]{}
\setkeys{Gin}{width=\maxwidth,height=\maxheight,keepaspectratio}
% Set default figure placement to htbp
\makeatletter
\def\fps@figure{htbp}
\makeatother

\usepackage{xcolor}
\usepackage{titlesec}
\AtBeginDocument{\renewcommand{\chaptername}{Capítulo}\renewcommand{\figurename}{Figura}}
\titleformat{\section}{\Large\bfseries\sffamily}{\thesection}{1em}{}
\titleformat{\subsection}{\large\bfseries\sffamily}{\thesubsection}{1em}{}
\titlespacing*{\section}{0pt}{1.5em}{0.75em}
\titlespacing*{\subsection}{0pt}{1.25em}{0.5em}
\makeatletter
\@ifpackageloaded{bookmark}{}{\usepackage{bookmark}}
\makeatother
\makeatletter
\@ifpackageloaded{caption}{}{\usepackage{caption}}
\AtBeginDocument{%
\ifdefined\contentsname
  \renewcommand*\contentsname{Table of contents}
\else
  \newcommand\contentsname{Table of contents}
\fi
\ifdefined\listfigurename
  \renewcommand*\listfigurename{List of Figures}
\else
  \newcommand\listfigurename{List of Figures}
\fi
\ifdefined\listtablename
  \renewcommand*\listtablename{List of Tables}
\else
  \newcommand\listtablename{List of Tables}
\fi
\ifdefined\figurename
  \renewcommand*\figurename{Figure}
\else
  \newcommand\figurename{Figure}
\fi
\ifdefined\tablename
  \renewcommand*\tablename{Table}
\else
  \newcommand\tablename{Table}
\fi
}
\@ifpackageloaded{float}{}{\usepackage{float}}
\floatstyle{ruled}
\@ifundefined{c@chapter}{\newfloat{codelisting}{h}{lop}}{\newfloat{codelisting}{h}{lop}[chapter]}
\floatname{codelisting}{Listing}
\newcommand*\listoflistings{\listof{codelisting}{List of Listings}}
\makeatother
\makeatletter
\makeatother
\makeatletter
\@ifpackageloaded{caption}{}{\usepackage{caption}}
\@ifpackageloaded{subcaption}{}{\usepackage{subcaption}}
\makeatother

\ifLuaTeX
  \usepackage{selnolig}  % disable illegal ligatures
\fi
\usepackage{bookmark}

\IfFileExists{xurl.sty}{\usepackage{xurl}}{} % add URL line breaks if available
\urlstyle{same} % disable monospaced font for URLs
\hypersetup{
  pdftitle={Microeconomía Intermedia},
  pdfauthor={Eduardo Zambrano},
  colorlinks=true,
  linkcolor={ForestGreen},
  filecolor={Maroon},
  citecolor={Blue},
  urlcolor={ForestGreen},
  pdfcreator={LaTeX via pandoc}}


\title{Microeconomía Intermedia}
\usepackage{etoolbox}
\makeatletter
\providecommand{\subtitle}[1]{% add subtitle to \maketitle
  \apptocmd{\@title}{\par {\large #1 \par}}{}{}
}
\makeatother
\subtitle{Paso a paso}
\author{Eduardo Zambrano}
\date{2026-10-13}

\begin{document}
\maketitle

\renewcommand*\contentsname{Tabla de Contenidos}
{
\hypersetup{linkcolor=}
\setcounter{tocdepth}{2}
\tableofcontents
}

\setstretch{1.15}
\bookmarksetup{startatroot}

\chapter*{Preámbulo}\label{preuxe1mbulo}
\addcontentsline{toc}{chapter}{Preámbulo}

\markboth{Preámbulo}{Preámbulo}

Estas notas acompañan a la asignatura de
\href{https://cvnet.cpd.ua.es/Guia-Docente/GuiaDocente/Index?wlengua=es&wcodasi=22014&scaca=2026-27\#}{Microeconomía
Intermedia} de la \href{https://www.ua.es/}{Universidad de Alicante}.

Su propósito es que tengas a mano, resueltos de principio a fin, los
cálculos que aparecen una y otra vez en la asignatura. Cada capítulo
toma ejemplos concretos y los desarrolla paso a paso para que puedas
seguir la mecánica con lápiz y papel y comprobar tus propios resultados.

Las notas se centran en el \emph{cómo} y no en el \emph{porqué}: la
intuición económica y la interpretación de los resultados se trabajan en
clase y en el manual. Lo que encontrarás aquí son las cuentas.

\section*{Uso de inteligencia
artificial}\label{uso-de-inteligencia-artificial}
\addcontentsline{toc}{section}{Uso de inteligencia artificial}

\markright{Uso de inteligencia artificial}

Estas notas se han elaborado con la ayuda de Claude Code (versión
2.1.296), el agente de programación de Anthropic, con el modelo Claude
Opus 5.5 (\texttt{claude-opus-5-5}), en octubre de 2026. Bajo mi
dirección, la IA redactó el texto, realizó las derivaciones, escribió el
código que genera los gráficos y construyó este sitio web. Yo elegí los
ejemplos y los valores de los parámetros, especifiqué el contenido y la
estructura de cada capítulo y revisé el material.

--- Eduardo Zambrano

\bookmarksetup{startatroot}

\chapter{Paso a paso: funciones de costes a corto y a largo
plazo}\label{paso-a-paso-funciones-de-costes-a-corto-y-a-largo-plazo}

\section{Datos comunes}\label{datos-comunes}

Una empresa produce \(q\) unidades de producto con trabajo, \(L\), y
capital, \(K\):

\[q = f(L, K)\]

El precio del trabajo es \(w\) y el precio del capital es \(r\):

\[w = 4 \qquad\qquad r = 1\]

El coste de producir con \((L, K)\) es:

\[CT = wL + rK\]

A largo plazo (\(lp\)) la empresa elige \(L\) y \(K\). A corto plazo
(\(cp\)) el capital está fijo en \(\bar K\) y la empresa solo elige
\(L\):

\[\bar K = 8\]

Tres tecnologías:

\[q = L^{1/3}K^{2/3} \qquad\quad q = \tfrac13 L + \tfrac23 K \qquad\quad q = \min\left\{\tfrac13 L,\ \tfrac23 K\right\}\]

\section{\texorpdfstring{Cobb-Douglas:
\(q = L^{1/3}K^{2/3}\)}{Cobb-Douglas: q = L\^{}\{1/3\}K\^{}\{2/3\}}}\label{cobb-douglas-q-l13k23}

\subsection{Corto plazo: demanda condicionada de
trabajo}\label{corto-plazo-demanda-condicionada-de-trabajo}

Con \(\bar K = 8\):

\[\bar K^{2/3} = 8^{2/3} = 4\]

Sustituyendo en la función de producción:

\[q = L^{1/3}\,\bar K^{2/3} = 4\,L^{1/3}\]

Despejando \(L\):

\[L^{1/3} = \frac{q}{4} \;\;\Longrightarrow\;\; L(q) = \frac{q^3}{64}\]

\subsection{Corto plazo: función de
costes}\label{corto-plazo-funciuxf3n-de-costes}

\[CT_{cp}(q) = w\,L(q) + r\,\bar K = 4\cdot\frac{q^3}{64} + 1\cdot 8\]

Simplificando:

\[CT_{cp}(q) = \underbrace{\frac{q^3}{16}}_{CV}\; +\; \underbrace{8}_{CF}\]

\begin{center}
\includegraphics{costes_files/figure-pdf/cell-3-output-1.pdf}
\end{center}

\subsection{Corto plazo: coste marginal y coste
medio}\label{corto-plazo-coste-marginal-y-coste-medio}

Derivando \(CT_{cp}\):

\[CMg_{cp}(q) = \frac{dCT_{cp}}{dq} = \frac{3q^2}{16}\]

Dividiendo \(CT_{cp}\) entre \(q\):

\[CMe_{cp}(q) = \frac{CT_{cp}}{q} = \frac{q^2}{16} + \frac{8}{q}\]

Mínimo del \(CMe_{cp}\):

\[\frac{dCMe_{cp}}{dq} = \frac{q}{8} - \frac{8}{q^2} = 0 \;\;\Longrightarrow\;\; q^3 = 64 \;\;\Longrightarrow\;\; q = 4\]

Evaluando en \(q = 4\):

\[CMe_{cp}(4) = 1 + 2 = 3 \qquad\qquad CMg_{cp}(4) = \frac{3\cdot 16}{16} = 3\]

\begin{center}
\includegraphics{costes_files/figure-pdf/cell-4-output-1.pdf}
\end{center}

\subsection{Largo plazo: condición de
tangencia}\label{largo-plazo-condiciuxf3n-de-tangencia}

Productos marginales:

\[PMg_L = \tfrac13 L^{-2/3}K^{2/3} \qquad\qquad PMg_K = \tfrac23 L^{1/3}K^{-1/3}\]

Dividiendo:

\[RTS = \frac{PMg_L}{PMg_K} = \frac{K}{2L}\]

Igualando a \(w/r\):

\[\frac{K}{2L} = \frac{4}{1} \;\;\Longrightarrow\;\; K = 8L\]

\subsection{Largo plazo: demandas
condicionadas}\label{largo-plazo-demandas-condicionadas}

Sustituyendo \(K = 8L\) en la función de producción:

\[q = L^{1/3}(8L)^{2/3} = 4L\]

Despejando:

\[L(q) = \frac{q}{4} \qquad\qquad K(q) = 8\,L(q) = 2q\]

\subsection{Largo plazo: función de
costes}\label{largo-plazo-funciuxf3n-de-costes}

\[CT_{lp}(q) = w\,L(q) + r\,K(q) = 4\cdot\frac{q}{4} + 1\cdot 2q\]

Simplificando:

\[CT_{lp}(q) = 3q\]

Coste marginal y coste medio:

\[CMg_{lp}(q) = 3 \qquad\qquad CMe_{lp}(q) = 3\]

\subsection{Isocuanta e isocoste
tangente}\label{isocuanta-e-isocoste-tangente}

Para \(q = 4\):

\[L(4) = 1 \qquad K(4) = 8 \qquad CT_{lp}(4) = 12\]

\begin{center}
\includegraphics{costes_files/figure-pdf/cell-5-output-1.pdf}
\end{center}

\subsection{Funciones de costes a corto y a largo
plazo}\label{funciones-de-costes-a-corto-y-a-largo-plazo}

\begin{center}
\includegraphics{costes_files/figure-pdf/cell-6-output-1.pdf}
\end{center}

\subsection{Rendimientos a escala}\label{rendimientos-a-escala}

\[f(tL,\,tK) = (tL)^{1/3}(tK)^{2/3} = t^{1/3+2/3}\,L^{1/3}K^{2/3}\]

Por tanto:

\[f(tL,\,tK) = t\,f(L,K)\]

La tecnología presenta rendimientos constantes a escala.

\section{\texorpdfstring{Sustitutivos perfectos:
\(q = \tfrac13 L + \tfrac23 K\)}{Sustitutivos perfectos: q = \textbackslash tfrac13 L + \textbackslash tfrac23 K}}\label{sustitutivos-perfectos-q-tfrac13-l-tfrac23-k}

\subsection{Corto plazo: demanda condicionada de
trabajo}\label{corto-plazo-demanda-condicionada-de-trabajo-1}

Con \(\bar K = 8\):

\[q = \tfrac13 L + \tfrac23 \bar K = \tfrac13 L + \tfrac{16}{3}\]

Despejando \(L\):

\[L(q) = \begin{cases} 0 & q \le \tfrac{16}{3} \\[4pt] 3q - 16 & q > \tfrac{16}{3} \end{cases}\]

\subsection{Corto plazo: función de
costes}\label{corto-plazo-funciuxf3n-de-costes-1}

\[CT_{cp}(q) = w\,L(q) + r\,\bar K = 4\,L(q) + 1\cdot 8\]

Sustituyendo \(L(q)\):

\[CT_{cp}(q) = \begin{cases} 8 & q \le \tfrac{16}{3} \\[4pt] 12q - 56 & q > \tfrac{16}{3} \end{cases}\]

\begin{center}
\includegraphics{costes_files/figure-pdf/cell-7-output-1.pdf}
\end{center}

\subsection{Corto plazo: coste marginal y coste
medio}\label{corto-plazo-coste-marginal-y-coste-medio-1}

Derivando \(CT_{cp}\):

\[CMg_{cp}(q) = \begin{cases} 0 & q < \tfrac{16}{3} \\[4pt] 12 & q > \tfrac{16}{3} \end{cases}\]

Dividiendo \(CT_{cp}\) entre \(q\):

\[CMe_{cp}(q) = \begin{cases} \dfrac{8}{q} & q \le \tfrac{16}{3} \\[8pt] 12 - \dfrac{56}{q} & q > \tfrac{16}{3} \end{cases}\]

Mínimo del \(CMe_{cp}\), en \(q = \tfrac{16}{3}\):

\[CMe_{cp}\left(\tfrac{16}{3}\right) = \frac{8}{16/3} = \frac32\]

\begin{center}
\includegraphics{costes_files/figure-pdf/cell-8-output-1.pdf}
\end{center}

\subsection{\texorpdfstring{Largo plazo: comparar \(RTS\) y
\(w/r\)}{Largo plazo: comparar RTS y w/r}}\label{largo-plazo-comparar-rts-y-wr}

Productos marginales:

\[PMg_L = \tfrac13 \qquad\qquad PMg_K = \tfrac23\]

Dividiendo:

\[RTS = \frac{1/3}{2/3} = \frac12 \qquad\qquad \frac{w}{r} = 4\]

Comparando:

\[RTS < \frac{w}{r} \;\;\Longrightarrow\;\; L = 0\]

\subsection{Largo plazo: demandas
condicionadas}\label{largo-plazo-demandas-condicionadas-1}

Sustituyendo \(L = 0\) en la función de producción:

\[q = \tfrac23 K\]

Despejando:

\[L(q) = 0 \qquad\qquad K(q) = \tfrac32 q\]

\subsection{Largo plazo: función de
costes}\label{largo-plazo-funciuxf3n-de-costes-1}

\[CT_{lp}(q) = w\,L(q) + r\,K(q) = 4\cdot 0 + 1\cdot\tfrac32 q\]

Simplificando:

\[CT_{lp}(q) = \tfrac32 q\]

Coste marginal y coste medio:

\[CMg_{lp}(q) = \tfrac32 \qquad\qquad CMe_{lp}(q) = \tfrac32\]

\subsection{Isocuanta e isocoste}\label{isocuanta-e-isocoste}

Para \(q = 4\):

\[L(4) = 0 \qquad K(4) = 6 \qquad CT_{lp}(4) = 6\]

\begin{center}
\includegraphics{costes_files/figure-pdf/cell-9-output-1.pdf}
\end{center}

\subsection{Funciones de costes a corto y a largo
plazo}\label{funciones-de-costes-a-corto-y-a-largo-plazo-1}

\begin{center}
\includegraphics{costes_files/figure-pdf/cell-10-output-1.pdf}
\end{center}

\subsection{Rendimientos a escala}\label{rendimientos-a-escala-1}

\[f(tL,\,tK) = \tfrac13 (tL) + \tfrac23 (tK) = t\left(\tfrac13 L + \tfrac23 K\right)\]

Por tanto:

\[f(tL,\,tK) = t\,f(L,K)\]

La tecnología presenta rendimientos constantes a escala.

\section{\texorpdfstring{Complementarios perfectos:
\(q = \min\left\{\tfrac13 L,\ \tfrac23 K\right\}\)}{Complementarios perfectos: q = \textbackslash min\textbackslash left\textbackslash\{\textbackslash tfrac13 L,\textbackslash{} \textbackslash tfrac23 K\textbackslash right\textbackslash\}}}\label{complementarios-perfectos-q-minlefttfrac13-l-tfrac23-kright}

\subsection{Corto plazo: demanda condicionada de
trabajo}\label{corto-plazo-demanda-condicionada-de-trabajo-2}

Con \(\bar K = 8\):

\[q = \min\left\{\tfrac13 L,\ \tfrac23 \bar K\right\} = \min\left\{\tfrac13 L,\ \tfrac{16}{3}\right\} \;\;\Longrightarrow\;\; q \le \tfrac{16}{3}\]

Despejando \(L\):

\[L(q) = 3q \qquad \left(q \le \tfrac{16}{3}\right)\]

\subsection{Corto plazo: función de
costes}\label{corto-plazo-funciuxf3n-de-costes-2}

\[CT_{cp}(q) = w\,L(q) + r\,\bar K = 4\cdot 3q + 1\cdot 8\]

Simplificando:

\[CT_{cp}(q) = \underbrace{12q}_{CV}\; +\; \underbrace{8}_{CF} \qquad \left(q \le \tfrac{16}{3}\right)\]

\begin{center}
\includegraphics{costes_files/figure-pdf/cell-11-output-1.pdf}
\end{center}

\subsection{Corto plazo: coste marginal y coste
medio}\label{corto-plazo-coste-marginal-y-coste-medio-2}

Derivando \(CT_{cp}\):

\[CMg_{cp}(q) = 12 \qquad \left(q \le \tfrac{16}{3}\right)\]

Dividiendo \(CT_{cp}\) entre \(q\):

\[CMe_{cp}(q) = 12 + \frac{8}{q} \qquad \left(q \le \tfrac{16}{3}\right)\]

Mínimo del \(CMe_{cp}\), en \(q = \tfrac{16}{3}\):

\[CMe_{cp}\left(\tfrac{16}{3}\right) = 12 + \frac{8}{16/3} = \frac{27}{2}\]

\begin{center}
\includegraphics{costes_files/figure-pdf/cell-12-output-1.pdf}
\end{center}

\subsection{Largo plazo: demandas
condicionadas}\label{largo-plazo-demandas-condicionadas-2}

Sin desperdiciar ningún factor:

\[\tfrac13 L = \tfrac23 K = q\]

Despejando:

\[L(q) = 3q \qquad\qquad K(q) = \tfrac32 q\]

\subsection{Largo plazo: función de
costes}\label{largo-plazo-funciuxf3n-de-costes-2}

\[CT_{lp}(q) = w\,L(q) + r\,K(q) = 4\cdot 3q + 1\cdot \tfrac32 q\]

Simplificando:

\[CT_{lp}(q) = \tfrac{27}{2}\, q\]

Coste marginal y coste medio:

\[CMg_{lp}(q) = \tfrac{27}{2} \qquad\qquad CMe_{lp}(q) = \tfrac{27}{2}\]

\subsection{Isocuanta e isocoste}\label{isocuanta-e-isocoste-1}

Para \(q = 2\):

\[L(2) = 6 \qquad K(2) = 3 \qquad CT_{lp}(2) = 27\]

\begin{center}
\includegraphics{costes_files/figure-pdf/cell-13-output-1.pdf}
\end{center}

\subsection{Funciones de costes a corto y a largo
plazo}\label{funciones-de-costes-a-corto-y-a-largo-plazo-2}

\begin{center}
\includegraphics{costes_files/figure-pdf/cell-14-output-1.pdf}
\end{center}

\subsection{Rendimientos a escala}\label{rendimientos-a-escala-2}

\[f(tL,\,tK) = \min\left\{\tfrac13 tL,\ \tfrac23 tK\right\} = t\,\min\left\{\tfrac13 L,\ \tfrac23 K\right\}\]

Por tanto:

\[f(tL,\,tK) = t\,f(L,K)\]

La tecnología presenta rendimientos constantes a escala.




\end{document}
